problem:  
Let \((M,\mathcal{D},g_{\mathcal{D}})\) be a compact, equiregular sub‑Riemannian manifold of step \(2\) equipped with its Popp measure \(\mu_{\mathrm{Popp}}\).  
For two smooth complements \(\mathcal{V}_1,\mathcal{V}_2\) with uniformly bounded O’Neill tensors, define Riemannian approximations  
\[
g_\lambda^{(i)} = g_{\mathcal{D}}\oplus \lambda^2 g_{\mathcal{V}_i},\qquad \lambda>0,\; i=1,2.
\]
Let \(\mathrm{Ric}_{g_\lambda^{(i)}}\) be the Ricci tensors of these metrics, and for any fixed horizontal field \(X\in\Gamma(\mathcal{D})\), set
\[
\mathsf{R}_\lambda^{(i)}(X)=\lambda^{-2}\,\mathrm{Ric}_{g_\lambda^{(i)}}(X,X).
\]

**Question:**  
Does one always have
\[
\lim_{\lambda\to\infty}
\big(
\mathsf{R}_\lambda^{(1)}(X)
-
\mathsf{R}_\lambda^{(2)}(X)
\big)
=0
\]
in the sense of distributions on \(M\)?
Equivalently, is the renormalized horizontal Ricci tensor \(\lambda^{-2}\mathrm{Ric}_{g_\lambda}\big|_{\mathcal{D}\times\mathcal{D}}\) **asymptotically independent** of the chosen complement when O’Neill tensors remain uniformly bounded?

If this fails in general, determine the precise geometric or analytic obstructions (e.g. non‑fatness of \(\mathcal{D}\), non‑integrability of higher commutators, or properties of torsion terms) and characterize the classes of sub‑Riemannian manifolds for which such complement‑independence holds.

---

Why is it a "good" problem:  
This question isolates one of the most elusive structural aspects lurking behind attempts to unify Riemannian and sub‑Riemannian curvature concepts: **complement dependence**. The proposed limit formulation removes ambiguities about arbitrary correction factors and asks a crisp asymptotic equality test—the simplest analytic criterion for whether Riemannian Ricci limits can define a geometric invariant intrinsic to the sub‑Riemannian manifold itself.  

Mathematically, the problem is fully well‑defined: each curvature term is computed from a smooth Riemannian metric, the scaling \(\lambda^{-2}\) corresponds to the known blow‑up of vertical directions in step‑2 geometries, and convergence is requested in a precise distributional sense. The question remains open even for contact structures, where the O’Neill tensors have rich algebraic structure but no established asymptotic invariance.  

Conceptually, resolving it would reveal whether sub‑Riemannian curvature can emerge as a limit independent of auxiliary choices—a foundational step toward an intrinsic “unified curvature” theory. Its concise form hides a profound interplay of collapse analysis, differential geometry, and control theory, exemplifying the depth, cross‑disciplinary reach, and elegant simplicity of a *good* research‑level mathematics problem.